\scanpage{11}
\mspage{- 101 -}
% archive sheet number: 10
% lines 1
$X$\handwritten{\,}\uline{des composantes irréductibles de} $Y$.
% slanted
\margin{\handwritten{codim. en \uline{un} point ...}}

% lines 2–3
\uline{Proposition 1} .- (i) \uline{Si} $Y$ \uline{est une partie d'un espace topologique} $X$ $\dim(Y)\handwritten{\leq}\dim(X)$.

% lines 4–5
(ii) \uline{Si un espace topologique} $X$ \uline{est réunion finie de parties fermé\handwritten{es}} $X_i$, \uline{alors} $\dim(X)=\sup_i(\dim(X_i))$.

% lines 6
\deleted{(iii) \uline{Pour toute \unsure{application continue}} \illegible{4} $\unsure{f(X)}$ \unsure{on a} \unsure{$(i)$}.}

% lines 7–10
Pour toute partie irréductible $Z$ fermée dans $Y$, l'adhérence $\overline Z$ de $Z$ dans $X$ est irréductible ; d'où (i). Si $X=\handwritten{\bigcup_i}X_i$, les $X_i$ étant f\handwritten{er}més, toute chaîne de parties fermées irréductibles de $X$ est nécessairement contenue dans l'un des $X_i$, ce qui prouve (ii).

% lines 10–12
\deleted{\unsure{Pour} (iii) \illegible{7} $\unsure{Y_x}$ \unsure{est une} \illegible{3} \unsure{application de} $\unsure{X}$ \unsure{(i)}.}

% lines 12–17
On notera que si $f$ est une application continue de $X$ dans un espace topologique $Y$, on peut avoir $\dim(f(X))\handwritten{>}\dim(X)$ ; on en a un exemple en prenant pour $X$ un espace discret à 2 éléments $a,b$, pour $Y$ l'ensemble $\handwritten{\{}a,b\handwritten{\}}$ muni de la topo\handwritten{lo}gie dans laquelle $\varnothing$, $\handwritten{\{}a\handwritten{\}}$ et $\handwritten{\{}a,b\handwritten{\}}$ sont les seuls ensembles fermé\handwritten{s} si $f$ est l'application identique $X\handwritten{\to}Y$, $\dim Y=1$ et $\dim X=0$.

% lines 18–20
\uline{Corollaire .- Soient} $x\handwritten{\in}X$, $U$ \uline{un voisinage de} $\handwritten{x}$, $Y_i$ $(1\handwritten{\leq}i\handwritten{\leq}n)$ \uline{des parties fermées de} $U$ \uline{telles que} $x\handwritten{\in}Y_i$ \uline{pour} $1\handwritten{\leq}i\handwritten{\leq}n$ \uline{et que} $U$ \uline{soit l\handwritten{a} réunion des} $Y_i$. \uline{Alors} $\dim_x X=\sup_i(\dim_x(Y_i))$.

% lines 21–25
En effet, d'après la prop.1, $\dim_x(X)=\inf_V(\sup_i(\dim(Y_i\handwritten{\cap}V)))$ lorsqu\handwritten{e} $V$ parcourt l'ensemble des voisinages ouverts de $x$ contenus dans $U$. Le corollaire est évident si $\sup_i(\dim_x(Y_i))=+\handwritten{\infty}$ ; sinon, il y a un voisinage $V_0$ de $x$ tel que \deleted{$\unsure{\sup_i|X_i}$\illegible{1}} $\dim(Y_i\handwritten{\cap}V)=\dim_x(Y_i)$ pour tout $\handwritten{i}$ et tout $V\handwritten{\subset}V_0$.

% lines 26–27
\uline{Proposition 2} .- (i) \uline{Soit} $\handwritten{\Phi}$ \uline{l'ensemble des parties fermées irrédu\handwritten{cti}bles de} $X$. \uline{On a} $\dim(X)=\sup_{Y\handwritten{\in\phi}}(\operatorname{codim}_X(Y))$.

% lines 28
(ii) \uline{On a} $\dim(X)=\sup_{x\handwritten{\in}X}\dim_x(X)$.

% lines 29
(i) découle aussitôt des déf.1 et 2 . \deleted{\unsure{D'autre part il existe dans}}
